%% kapa: exam=1 round=26 year=2015 period=1 form=A subject=civil_law subject_ko=민법 seq=33 number=33
%% source: Q-Net 감정평가사 기출문제 (cst00309) artlSeq=5138001 file=extracted/1008/5138001/제26회_감정평가사 1차_기출/제26회_감정평가사 1차_기출/1교시_감평_A형.pdf
%% conversion: kordoc 4.24.51 -> tools/gen1.py (auto); review=proofread-sonnet; flags=none
\begin{question}{33}
부동산 점유취득시효에 관한 설명으로 옳지 않은 것은? (다툼이 있으면 판례에 따름)
\begin{choices}
  \item 점유취득시효가 완성된 경우 점유자가 시효기간 중에 수취한 과실은 소유자에게 반환할 필요가 없다.
  \item 점유자의 점유가 불법이라고 주장하는 소유자로부터 이의를 받은 사실이 있다 하더라도 그러한 사실만으로 곧 평온ㆍ공연한 점유가 부정되지 않는다.
  \item 취득시효 진행 중에 소유자가 소유권을 제3자에게 양도하고 등기를 이전한 후 시효가 완성된 경우, 점유자는 양수인에게 시효 완성을 이유로 소유권이전등기를 청구할 수 있다.
  \item 종중 부동산이 종중 대표자에게 적법하게 명의신탁되었는데 그 부동산에 대해 제3자의 점유에 의한 취득시효가 완성된 후 제3자 명의의 등기 전에 명의신탁이 해지되어 등기명의가 종중에게 이전된 경우, 특별한 사정이 없는 한 점유자는 종중에 대해 시효 완성을 주장할 수 있다.
  \item 점유취득시효가 완성된 경우, 점유자는 등기 없이는 그 부동산의 소유권을 주장할 수 없다.
\end{choices}
\end{question}
\answer{4}
